\documentclass[11pt,a4paper]{article}
\usepackage[margin=1in]{geometry}
\usepackage{graphicx}
\usepackage{amsmath,amssymb}
\usepackage{booktabs}
\usepackage{caption}
\usepackage[hidelinks]{hyperref}
\graphicspath{{../results/fe55_mc/}{../results/scan_waveform_halfpitch/}}

\newcommand{\Fe}{\textsuperscript{55}Fe}
\newcommand{\co}{CO\textsubscript{2}}

\title{A Data-Driven Monte Carlo of \Fe{} Signals\\ in the Thin Gap Chamber}
\author{TGC Garfield\texttt{++} simulation study}
\date{\today}

\begin{document}
\maketitle

\begin{abstract}
\noindent
This note describes a light, \emph{data-driven} Monte Carlo that synthesizes full \Fe{} detector
signals in a Thin Gap Chamber (TGC) by superposing the \emph{measured} single-electron response over a
realistic cloud of primary electrons, with no further Garfield\texttt{++} microsimulation. Each primary
electron is given its own drift delay and avalanche gain, drawn from position-resolved maps obtained in
a prior scan, and stamps the measured single-electron waveform; the sum is the event signal. The model
reproduces the characteristic \Fe{} charge spectrum (the 5.9~keV photopeak and the Ar escape peak at
half the charge) from first principles, and is then used to study the pulse shape, the energy
resolution, and their dependence on the conversion position and on the X-ray collimation. The tool is
fast (a full 50-point position grid plus a $2\times10^{4}$-photon exposure runs in about one minute on a
laptop) and is implemented in \texttt{tools/fe55\_mc.py}.
\end{abstract}

\tableofcontents

\section{Introduction}
The TGC is a multiwire proportional chamber run in a quasi-saturated mode: a plane of
50~$\mu$m-diameter anode wires at 1.8~mm pitch sits midway between two
grounded cathode planes 1.4~mm away on each side, filled with Ar/\co{} 70:30 at
1900~V. An \Fe{} source emits Mn~K$\alpha$ X-rays at 5.9~keV, which photo-convert in the
gas and drift to the nearest wire, where they avalanche.

Characterizing the \emph{signal} — its shape, amplitude, and the resulting energy resolution — normally
requires an event-by-event microscopic avalanche simulation, which is expensive. The key observation
that makes a fast model possible is that the single-electron response \emph{factorizes}: a prior scan
(Sec.~\ref{sec:method}) established that the peak-aligned, peak-normalized single-electron anode pulse
shape $S(t')$ is \emph{position-independent}, while the electron drift delay and the avalanche gain are
strongly position-dependent and can be tabulated. A real event is a linear superposition of
single-electron responses, so the full signal can be assembled by summing the measured shape over a
realistic set of primaries without any new field or avalanche computation.

\section{Method}
\label{sec:method}

\subsection{Factorization and measured inputs}
The simulation consumes three products of the position scan, all defined on the same half-pitch grid of
transverse distance-from-wire $d$ and conversion depth $y$ (folded into one wire cell by the glide
symmetry of the wire plane):
\begin{itemize}
  \item the canonical single-electron anode shape $S(t')$ (position-independent, pooled over all scan
        positions; peak-normalized with the current spike at $t'=0$);
  \item the drift-delay map $\bar\tau(d,y)$ and its event-to-event jitter $\sigma_\tau(d,y)$ (the latter
        carries the single-electron longitudinal diffusion);
  \item the per-position avalanche-gain distribution (the full Polya, including the finite probability
        of zero gain from attachment), measured as the per-event avalanche size.
\end{itemize}
Figure~\ref{fig:shape} shows the measured single-electron response. The peak-aligned shapes overlay
across positions (upper panels), confirming the position independence of $S(t')$ that the model relies
on; the lower panels show that the electron drift time grows smoothly with depth and distance from the
wire.

\begin{figure}[htbp]
  \centering
  \includegraphics[width=\textwidth]{waveform_shapes.png}
  \caption{The \textbf{measured single-electron response} used as the model input (ion drift on).
  Upper left: peak-aligned, peak-normalized anode pulses at the four corner positions, overlaying to a
  common shape. Upper right: a zoom on the first 10~ns (spike plus the start of the ion
  tail). Lower left: the ion-tail height across the cell. Lower right: the electron drift time to the
  wire. The position-independence of the upper-panel shape is what licenses the superposition model.}
  \label{fig:shape}
\end{figure}

\subsection{Signal synthesis}
For an event with $N$ primary electrons, the induced anode current is built as
\begin{equation}
  W(t) \;=\; \sum_{i=1}^{N} g_i \, S\!\left(t-\tau_i\right),
  \qquad
  g_i \sim \mathrm{Gain}(d_i,y_i),
  \qquad
  \tau_i \sim \mathcal{N}\!\big(\bar\tau,\sigma_\tau\big)(d_i,y_i),
  \label{eq:sum}
\end{equation}
\noindent where the symbols are, term by term:
\begin{description}
  \item[$W(t)$] the event's induced anode current --- the signal --- as a function of time $t$ measured
        from the photon conversion. Its units are arbitrary (gain-electrons $\times$ the peak-normalized
        shape); only its shape and relative amplitude are used.
  \item[$N$, $i$] the number of primary ionization electrons the photon liberates, and the index
        $i=1\dots N$ running over them. $N\approx E_\mathrm{dep}/W$ with a Fano fluctuation
        (Sec.~\ref{sec:primary}): $N\approx227$ for a fully absorbed 5.9~keV photon and $N\approx113$
        for an Ar-escape event.
  \item[$g_i$] the avalanche gain of primary $i$ --- the number of electrons its avalanche collects ---
        drawn \emph{independently} for each primary from the measured single-electron gain distribution
        $\mathrm{Gain}(d_i,y_i)$ at that primary's position. That distribution is the per-cell Polya
        measured by Garfield, including the finite probability of \emph{zero} gain (attachment of the
        primary before it multiplies).
  \item[$S(t')$] the measured, position-independent, peak-aligned and peak-normalized single-electron
        anode pulse shape (Fig.~\ref{fig:shape}), with its current spike at $t'=0$. The term
        $g_i\,S(t-\tau_i)$ is therefore primary $i$'s contribution: the universal single-electron pulse,
        scaled by that primary's gain $g_i$ and shifted to its arrival time $\tau_i$.
  \item[$\tau_i$] the drift time of primary $i$ from its creation point to the wire, sampled from a
        Gaussian $\mathcal{N}(\bar\tau,\sigma_\tau)$ with the measured mean $\bar\tau(d_i,y_i)$ and width
        $\sigma_\tau(d_i,y_i)$; $\bar\tau$ is the drift-time map and $\sigma_\tau$ its event-to-event
        RMS, i.e.\ the single-electron longitudinal-diffusion jitter.
  \item[$(d_i,y_i)$] the primary's position --- $d_i$ the transverse distance to the nearest wire and
        $y_i$ the conversion depth --- folded by the glide symmetry of the wire plane into the single
        measured half-pitch $\times$ depth cell on which the maps live. $\bar\tau$ and $\sigma_\tau$ are
        bilinearly interpolated there, while $\mathrm{Gain}$ is taken from the nearest measured cell.
\end{description}
The sum is accumulated efficiently as one convolution of the per-primary charge-arrival profile
$\sum_i g_i\,\delta(t-\tau_i)$ with $S$. Equation~\eqref{eq:sum} is exactly the ``unfolding'' of the
production simulation, which instead fires one representative avalanche and scales the result by
$N=\mathrm{round}(E/W)$; here the $N$ primaries are independent, each with its own position, arrival time
and gain --- which is what produces the realistic pulse shape and its event-to-event fluctuations.

\subsection{\Fe{} primary-charge generation}
\label{sec:primary}
The primary ionization is built microscopically, cluster by cluster, so that both the number of
primaries $N$ and their spatial pattern fluctuate realistically. This is the step that supplies the set
of primaries $\{(d_i,y_i)\}$ --- and hence $N$ --- consumed by Eq.~\eqref{eq:sum}.

\paragraph{Photoabsorption.} An \Fe{} source emits mainly Mn~K$\alpha$ at $E_\gamma=5.895$~keV. Above the
argon K-edge the photon is absorbed almost entirely on the Ar~K shell (binding energy $E_K=3.206$~keV;
the \co{} and the argon L shells are negligible absorbers at this energy), ejecting a photoelectron of
energy
\begin{equation*}
  E_\mathrm{pe} = E_\gamma - E_K \approx 2.69~\text{keV}.
\end{equation*}
The remaining $E_K$ stays momentarily as a K-shell vacancy, released as described below.

\paragraph{Photoelectron track.} The photoelectron is emitted with the dipole angular distribution
$\mathrm{d}N/\mathrm{d}\Omega \propto \sin^2\theta\,(1+\beta\cos\theta)$ relative to the incident-photon
direction ($\beta=v/c\approx0.10$ at 2.69~keV, so the pattern is nearly $\sin^2\theta$ with a weak
forward tilt) and ranges out over a short, tortuous track. Its practical range is parametrised as
$R(E)=R_0\,(E/\text{keV})^{1.7}$ with $R_0\approx30~\mu$m, giving $R(2.69~\text{keV})\approx0.16$~mm at
1~atm. The $n_\mathrm{pe}=E_\mathrm{pe}/W\approx103$ ionization electrons it creates ($W=26$~eV) are
placed uniformly along the emission segment, each displaced by a Gaussian transverse straggling of width
$\sigma_\mathrm{str}=f\,R$ ($f\approx0.3$). This gives the directional, finite-size charge blob of a real
soft-X-ray conversion rather than a point.

\paragraph{K-shell relaxation.} The Ar~K vacancy fills in one of two ways, chosen with the
K-fluorescence yield $\omega_K\approx0.12$:
\begin{itemize}
  \item \emph{Auger} (probability $1-\omega_K\approx0.88$): the $E_K=3.206$~keV is released as an Auger
        electron and its cascade, deposited locally as a second short track at the conversion point. The
        event then carries the full 5.9~keV, giving $N=E_\gamma/W\approx227$ primaries --- the
        \textbf{main (photopeak)} events.
  \item \emph{Fluorescence} (probability $\omega_K\approx0.12$): an Ar~K$\alpha$ photon of
        $E_{\mathrm{K}\alpha}=2.957$~keV is emitted isotropically and travels an exponentially
        distributed distance (attenuation length $\approx2$~cm at this energy). Because the gap is only
        $\sim$2.8~mm it usually \emph{escapes} the active volume: then only
        $E_\mathrm{pe}+E_L\approx2.94$~keV is deposited --- where $E_L=E_K-E_{\mathrm{K}\alpha}\approx
        0.249$~keV is the residual L-shell Auger left at the origin --- giving $N\approx113$, the
        \textbf{Ar-escape} events at half the charge. If the fluorescence photon is instead re-absorbed
        in the gas it lays down $E_{\mathrm{K}\alpha}$ as a displaced satellite cluster and the event
        returns to the photopeak.
\end{itemize}

\paragraph{Fano fluctuation.} The electron count of every energetic deposit is sub-Poisson: each
$n=E/W$ is fluctuated as $n\sim\mathcal{N}\!\big(E/W,\ \sqrt{F\,E/W}\big)$ with a Fano factor
$F\approx0.2$, contributing an intrinsic $\sigma_E/E=\sqrt{F/N}$ per event (a few percent) on top of the
avalanche (Polya) spread.

\paragraph{Cloud, diffusion, and hand-off to Eq.~\eqref{eq:sum}.} The union of the photoelectron track,
the Auger / L-Auger blobs and any fluorescence satellite is the primary-electron cloud. As it drifts,
each electron additionally acquires a transverse Gaussian spread $\sigma_x=D\sqrt{y}$
($D\approx120~\mu$m$/\sqrt{\text{cm}}$) before its $(d_i,y_i)$ is folded onto the maps; the longitudinal
(arrival-time) diffusion is already carried by $\sigma_\tau$. Every primary then contributes
$g_i\,S(t-\tau_i)$ exactly as in Eq.~\eqref{eq:sum}.

\paragraph{Exposure geometries.} Two illuminations are modelled. In the \textbf{realistic full-cell}
exposure the conversion depth follows the exponential X-ray absorption through the gap --- with
$\lambda(5.9~\text{keV})\approx16$~cm~$\gg$~gap this is nearly uniform in depth --- while the transverse
position is uniform over one wire pitch. In the \textbf{collimated} exposure the beam passes through a
circular aperture, so the conversion footprint is a disk in the plane transverse to the beam, centred on
a wire or between wires; its chord-weighted (semicircle) transverse profile concentrates conversions
near the aperture centre.

All numerical inputs above ($W$, $F$, $\omega_K$, the line and binding energies, $R_0$ and the range
exponent, the straggling fraction $f$, the diffusion coefficient $D$, and the attenuation lengths) are
set in \texttt{config/fe55\_mc.json}.

\clearpage
\section{Results}

\subsection{Charge spectrum and validation}
Figure~\ref{fig:spectrum} shows the collected-charge spectrum $Q=\sum_i g_i$ over the realistic
exposure. The model reproduces the classic \Fe{}-in-argon double-peak structure with no tuning:
\begin{itemize}
  \item the \textbf{main-peak centroid} equals $N\langle g\rangle \approx 227\times31800 \approx
        7.2\times10^{6}$ gain-electrons;
  \item the \textbf{escape/main charge ratio} is $0.50$, matching the energy ratio
        $2.94/5.9=0.498$;
  \item the \textbf{escape fraction} is $\approx 9.5\%$, consistent with $\omega_K$ times the
        geometric escape probability of the thin gap.
\end{itemize}
The energy resolution of the uniform exposure is $\sigma_E/E\approx19\%$. This is broader than a
textbook proportional counter because it is dominated not by the Polya and Fano terms
($\sqrt{(\sigma_g/\langle g\rangle)^2 + F}\,/\sqrt{N}\approx7\%$) but by the \emph{gain
non-uniformity across the cell} sampled by a full-cell exposure — see Sec.~\ref{sec:res} and~\ref{sec:perpoint}.

\begin{figure}[htbp]
  \centering
  \includegraphics[width=0.85\textwidth]{fe55_spectrum.png}
  \caption{\textbf{\Fe{} charge spectrum} (realistic full-cell exposure, $2\times10^{4}$ photons). Grey:
  all events; blue: full-energy (main) events; red: Ar-escape events. The top axis is self-calibrated by
  setting the main-peak centroid to 5.9~keV.}
  \label{fig:spectrum}
\end{figure}

\subsection{Pulse shape versus conversion position}
Aligning each mean pulse to its own peak removes the drift-time offset and isolates the intrinsic shape.
Figure~\ref{fig:waveforms} (left) shows that the shape depends on the \emph{transverse} position — sharp
on-wire, broad in the mid-gap — but is essentially independent of depth (the shallow/deep curves
overlay within each transverse position). Figure~\ref{fig:sweep} makes this explicit: sweeping the
transverse position at fixed depth broadens the pulse monotonically from the wire to the mid-gap, while
sweeping depth at fixed transverse position leaves the shape nearly unchanged.

\begin{figure}[htbp]
  \centering
  \includegraphics[width=\textwidth]{fe55_waveforms.png}
  \caption{Left: mean \Fe{} pulse at the four corner positions, \textbf{peak-aligned}; on-wire pulses
  are sharp, mid-gap pulses broad, and the shape is set by the transverse position, not the depth.
  Right: realistic single-event pulses about their aligned mean.}
  \label{fig:waveforms}
\end{figure}

\begin{figure}[htbp]
  \centering
  \includegraphics[width=\textwidth]{fe55_sweep.png}
  \caption{Aligned pulse shapes swept along one grid axis. Left: transverse sweep at fixed depth (color
  = distance from wire) — the pulse broadens smoothly toward the mid-gap. Right: depth sweep at fixed
  transverse position — the shape is nearly depth-independent.}
  \label{fig:sweep}
\end{figure}

The timing and width observables across the full cell are mapped in Fig.~\ref{fig:grid} and the pulse
FWHM in Fig.~\ref{fig:fwhm}. The peak time (drift time) grows from $\sim$1~ns
(shallow, on-wire) to $\sim$24~ns (deep, mid-gap); the 10--90\% rise time spans
1.1--7.1~ns and the FWHM 1.5--14.6~ns, both governed almost
entirely by the distance from the wire.

\begin{figure}[htbp]
  \centering
  \includegraphics[width=\textwidth]{fe55_grid.png}
  \caption{Pulse observables across the full (depth $\times$ distance-from-wire) grid: peak time (drift
  timing), 10--90\% rise time, mean peak amplitude, and mean charge.}
  \label{fig:grid}
\end{figure}

\begin{figure}[htbp]
  \centering
  \includegraphics[width=0.72\textwidth]{fe55_fwhm.png}
  \caption{Pulse FWHM versus position: $\sim$1.5--3~ns across the near-wire region,
  rising to $\sim$10--15~ns toward the mid-gap. The gradient is almost purely
  transverse.}
  \label{fig:fwhm}
\end{figure}

\clearpage
\subsection{Single-event inspection}
Figure~\ref{fig:single} shows individual events with their per-primary microstructure. On-wire, the
$\sim$227 primary arrivals cluster into a narrow time window and sum to a sharp spike; in the mid-gap
they spread over several nanoseconds, producing a broad, visibly lumpy pulse. Escape events are
immediately recognizable by their $N\!\approx\!113$ primaries and roughly half the charge.

\begin{figure}[p]
  \centering
  \includegraphics[width=\textwidth]{fe55_single.png}
  \caption{Single \Fe{} events (rows: the four characteristic positions; columns: three events each).
  Blue: induced current $W(t)$ at absolute time (real amplitude). Orange: the discrete per-primary
  charge arrivals (gain versus arrival time) that build the pulse.}
  \label{fig:single}
\end{figure}

\subsection{Energy resolution and collimation}
\label{sec:res}
Because the gain varies across the cell, the measured energy resolution depends on which part of the
cell the source illuminates. Figure~\ref{fig:compare} compares the charge spectrum of the full-cell
exposure with collimated exposures (a 1.8~mm-diameter circular aperture) centered on a
wire and between wires. Collimating \emph{between wires} gives the best resolution, because the gain is
flat and saturated in the mid-gap; collimating \emph{on a wire} is the worst, because it samples the
steep near-wire gain gradient. The escape peak is unchanged, as expected for a transverse-independent
process. The numbers are collected in Table~\ref{tab:summary}.

\begin{figure}[htbp]
  \centering
  \includegraphics[width=0.85\textwidth]{fe55_spectrum_compare.png}
  \caption{\textbf{Charge spectrum versus collimator position} (density-normalized). Between-wires
  collimation gives the tightest main peak; on-wire the broadest; the full-cell exposure lies in
  between.}
  \label{fig:compare}
\end{figure}

\subsection{Pulse-height (peak-amplitude) spectra}
A fast current readout measures the pulse height rather than the integrated charge. Because the peak
amplitude depends on the pulse \emph{width} as well as the charge ($A\propto Q/\mathrm{width}$), and the
width varies strongly across the cell, the pulse-height spectrum is far more position-smeared than the
charge spectrum: $\sigma_E/E\approx38\%$ for the uniform exposure (Fig.~\ref{fig:amp}), versus $19\%$
for charge. Under collimation the on-wire aperture is now the \emph{tighter} one (sharp pulses track the
charge directly), the reverse of the charge case (Fig.~\ref{fig:ampcompare}). The practical conclusion
is that, for this chamber, the integrated charge is a substantially better energy estimator than the
peak pulse height.

\begin{figure}[htbp]
  \centering
  \includegraphics[width=0.85\textwidth]{fe55_amplitude.png}
  \caption{\textbf{Pulse-height spectrum} (uniform exposure). The photopeak is broad and structured
  because the peak amplitude folds in the position-dependent pulse width.}
  \label{fig:amp}
\end{figure}

\begin{figure}[htbp]
  \centering
  \includegraphics[width=0.85\textwidth]{fe55_amplitude_compare.png}
  \caption{\textbf{Pulse-height spectrum versus collimator position}. On-wire (sharp pulses) is higher
  and tighter; between-wires (broad pulses) is lower; opposite to the charge ordering.}
  \label{fig:ampcompare}
\end{figure}

\clearpage
\subsection{Per-position distributions}
\label{sec:perpoint}
Figures~\ref{fig:gridcharge} and~\ref{fig:gridamp} show the charge and pulse-height distributions at
every one of the 50 grid points. At a \emph{fixed} position each main peak is sharp (a few percent),
confirming that the broad full-cell spectra of Figs.~\ref{fig:spectrum} and~\ref{fig:amp} arise from
mixing positions, not from an intrinsically poor single-position resolution. The per-point charge is
highest on-wire at mid-depth ($\sim$10500k) and lowest at the on-wire-shallow ``notch''
($\sim$3240k), where a primary born inside the avalanche region has little path to multiply; the
peak amplitude peaks on-wire at mid-depth and falls monotonically toward the mid-gap.

\begin{figure}[p]
  \centering
  \includegraphics[width=\textwidth]{fe55_grid_charge.png}
  \caption{Per-point \textbf{charge} distributions (rows: depth; columns: distance from wire; shared
  x-axis). Blue: main; red: escape. The main-peak centroid ($10^3$ gain-electrons) is annotated in each
  panel.}
  \label{fig:gridcharge}
\end{figure}

\begin{figure}[p]
  \centering
  \includegraphics[width=\textwidth]{fe55_grid_amplitude.png}
  \caption{Per-point \textbf{pulse-height} distributions, same layout. The amplitude is highest on-wire
  at mid-depth and decreases toward the mid-gap.}
  \label{fig:gridamp}
\end{figure}

\clearpage
\section{Summary}
A fast, data-driven Monte Carlo reproduces \Fe{} signals in the TGC by superposing the measured
single-electron response over a realistic primary-electron cloud. It validates against the known \Fe{}
spectrum (photopeak plus Ar escape at half the charge, correct centroid and escape fraction) and yields
the following picture of the chamber's signal:
\begin{itemize}
  \item the \textbf{pulse shape} is set by the transverse distance from the wire (sharp on-wire, broad
        mid-gap) and is nearly independent of depth, which only sets the drift delay;
  \item the \textbf{energy resolution} from integrated charge is $\approx19\%$ for a full-cell
        exposure, improving to $\approx17\%$ when collimated between wires and degrading to $\approx20\%$
        on a wire; at a fixed position it is intrinsically a few percent, so the spread is dominated by
        gain non-uniformity across the cell;
  \item the \textbf{peak pulse height} is a markedly poorer energy estimator ($\approx38\%$) than the
        integrated charge, because it also folds in the position-dependent pulse width.
\end{itemize}

\begin{table}[htbp]
  \centering
  \caption{Main-peak centroid and resolution of the \Fe{} photopeak for the three exposures, for the
  integrated charge and the peak amplitude. Charge in $10^3$ gain-electrons, amplitude in $10^6$
  arbitrary current units.}
  \label{tab:summary}
  \begin{tabular}{lcccc}
    \toprule
    & \multicolumn{2}{c}{Charge} & \multicolumn{2}{c}{Peak amplitude} \\
    \cmidrule(lr){2-3}\cmidrule(lr){4-5}
    Exposure & centroid & $\sigma_E/E$ & centroid & $\sigma_E/E$ \\
    \midrule
    Full-cell uniform        & 7220 & 19\% & 1.96 & 38\% \\
    Collimated, on-wire      & 7380 & 20\% & 2.12 & 35\% \\
    Collimated, between wires & 7020 & 17\% & 1.80 & 38\% \\
    \bottomrule
  \end{tabular}
\end{table}

\appendix
\section{Reproducing the figures}
All figures are produced by \texttt{tools/fe55\_mc.py} (pure Python; the maps come from the committed
position-scan products) into \texttt{results/fe55\_mc/}:
{\footnotesize
\begin{verbatim}
# position maps: charge spectrum, pulse shapes, observable maps, per-point distributions
python3 tools/fe55_mc.py --scan-grid --seed 1
# collimator study: charge + pulse-height, on-wire vs between-wires
python3 tools/fe55_mc.py --mode realistic --collimator 1.8 \
        --collimator-center both --seed 1
# single-event inspection
python3 tools/fe55_mc.py --inspect --seed 1
\end{verbatim}
}
See Section~14 of the TGC simulation manual (\texttt{docs/manual.md}) for the model details and the full
list of command-line options.

\end{document}
